\section{System Overview}
\label{sec:system}

AutoTerm-SST uses a thin framework with swappable translation agents, while its
main system contribution is the adaptive terminology loop inside the live agent
(Figure~\ref{fig:architecture}). The framework serves the web UI, accepts REST
control requests, streams audio over WebSockets, tracks session lifecycle, and
routes each session by \texttt{agent\_type}; it does not own model prompts,
retrieval, batching, KV cache, or terminology resources. For chunk \(t\), the
agent retrieves from the current budgeted working set before decoding. The
generated translation may update that working set for a later chunk, no earlier
than \(t+1\), but
never changes the prompt already used for chunk \(t\).

\begin{figure*}[t]
\centering
\includegraphics[width=0.98\textwidth]{autoterm_architecture.pdf}
\caption{AutoTerm-SST on a stream that moves from NLP to medicine. Each
speech chunk is encoded by the RASST backbone, and at most 10 bilingual term
    pairs retrieved from the \emph{active slices} are injected into the
language-model prompt. Sliding-window translation context ranks bilingual
slice prototypes under a fixed term budget, so \emph{SRL} is served from the
NLP glossary early and \emph{rectal cancer} from the medicine glossary later.
The working set may hold several slices; the ellipsis denotes additional
registered glossaries.}
\label{fig:architecture}
\end{figure*}

\paragraph{Agent contract.}
Agents implement a small contract: \texttt{open\_session},
\texttt{submit\_audio}, \texttt{on\_control}, \texttt{close\_session},
\texttt{describe}, and \texttt{health}. The same UI and wire protocol support a
deterministic mock agent for development, a legacy InfiniSST scheduler agent
\citep{ouyang2025infinisst}, or the AutoTerm-SST live agent built on
RASST/Qwen3-Omni \citep{luo2026rasst}. The boundary is
backend-agnostic: the live demo uses in-process vLLM \citep{kwon2023efficient},
and the same protocol can serve external SGLang-compatible backends
\citep{zheng2024sglang} without changing the terminology memory, router, or
UI.

\paragraph{Streaming agent.}
For each audio chunk, the OmniAgent scores the accumulated target-text window
against offline bilingual slice prototypes. It selects the highest-similarity
slices that fit the active-term budget, then runs MaxSim retrieval across that
working set. The agent filters and jointly reranks candidates, caps prompt
references at 10, and lets
the PromptBuilder format the \texttt{term\_map} before calling the in-process
vLLM backend. After decoding, \(y_t\) is appended to the target history for a
later routing tick. The agent emits partial translations and
structured metadata, including retrieved terms, retrieval latency, active
slices, prompt-reference counts, and router decisions.

\paragraph{Evidence protocol.}
The default WebSocket output remains plain text for backward compatibility. A
JSON mode returns \texttt{\{type, text, meta\}}, where
\texttt{meta.references} drives the live evidence panel and
\texttt{meta.topic} and \texttt{meta.topic\_router} describe the adaptive
glossary state. This evidence protocol lets users see which terms were retrieved
and whether the router stayed, switched, or deferred.
